
@@M0@@
où @@M1@@ est une
immersion régulière, et @@M2@@ un morphisme lisse. Comme @@M3@@ est lisse, c'est un morphisme parfait (III 4.1). D'après VII 1.9, @@M4@@ est également un morphisme parfait. Comme le composé de deux morphismes parfaits est parfait (III 4.5), on en déduit la proposition.

\textbf{Proposition 1.8.} \uline{Soient} @@M5@@ \uline{un morphisme d'intersection complète}, @@M6@@ \uline{un point de} @@M7@@, \uline{et une factorisation} :
@@M8@@
\uline{où} @@M9@@ \uline{est un voisinage ouvert de} @@M10@@, @@M11@@ \uline{une immersion régulière}, \uline{et} @@M12@@ \uline{un morphisme lisse}. \uline{On note} @@M13@@ \uline{l'idéal de} @@M14@@ \uline{définissant} @@M15@@. \uline{Alors l'entier}
@@M16@@
\uline{est indépendant de la factorisation choisie de} @@M17@@.

Supposons que la restriction de @@M18@@ à @@M19@@ se factorise par deux @@M20@@-schémas lisses @@M21@@ et @@M22@@, et posons @@M23@@. Alors @@M24@@ est lisse sur @@M25@@, et aussi sur @@M26@@ et @@M27@@. De plus, les immersions @@M28@@ et @@M29@@ définissent une immersion @@M30@@, régulière d'après 1.2. On peut donc se ramener à comparer les entiers @@M31@@ relatifs à deux factorisations de @@M32@@ par des @@M33@@-schémas lisses @@M34@@ et @@M35@@ tels qu'il existe un diagramme commutatif au-dessus de @@M36@@ :
@@M37@@
